\documentclass[../DoAn.tex]{subfiles}
\begin{document}


Chương 3 tập trung trình bày cơ sở lý thuyết và các công nghệ cốt lõi được sử dụng để giải quyết các bài toán kỹ thuật đã nêu trong Chương 2. Nội dung chương phân tích lý do lựa chọn từng thành phần trong hệ sinh thái công nghệ của đồ án, đối chiếu với một số giải pháp thay thế tương đương và làm rõ mức độ phù hợp của bộ giải pháp MERN Stack kết hợp Redis, Socket.IO, các cổng thanh toán trực tuyến và cơ chế khuyến nghị sản phẩm cá nhân hóa.

\section{Nền tảng cơ sở dữ liệu và Cơ chế giữ kho nguyên tử}
\label{section:3.1}

\subsection{Cơ sở lý thuyết về nhất quán dữ liệu và tranh chấp tài nguyên tải cao}
Trong các hệ thống thương mại điện tử quy mô lớn, tính nhất quán dữ liệu ở tầng lưu trữ là yếu tố sống còn quyết định sự thành bại của doanh nghiệp. Vấn đề tranh chấp dữ liệu khi có hàng nghìn người cùng thực hiện mua một sản phẩm tại cùng một thời điểm (kịch bản Flash Sale) thường phát sinh lỗi nghiêm trọng về bán quá mức tồn kho thực tế (\textit{Overselling}).

Về mặt lý thuyết, lỗi này xảy ra do mô hình xử lý giao dịch cổ điển áp dụng kỹ thuật đọc - chỉnh sửa - ghi (\textit{Read-Modify-Write}). Khi hai luồng yêu cầu đồng thời kiểm tra số lượng tồn kho của một sản phẩm (ví dụ còn $1$ sản phẩm), cả hai đều nhận được kết quả còn hàng và cùng tiến hành thực hiện trừ kho và tạo đơn, dẫn đến số lượng tồn kho thực tế bị âm ($-1$). Để giải quyết bài toán này từ Chương 2, hệ thống bắt buộc phải thực hiện thao tác kiểm tra điều kiện và cập nhật số lượng tồn kho một cách nguyên tử (\textit{Atomic Database Reservation}) kết hợp với cơ chế chống trùng lặp giao dịch bằng khóa \texttt{idempotencyKey} ở tầng ứng dụng \cite{nosql-performance}.

\subsection{Cơ sở dữ liệu tài liệu MongoDB và thư viện Mongoose ORM}
MongoDB là một cơ sở dữ liệu phi quan hệ (NoSQL) hướng tài liệu (\textit{Document-oriented}) lưu trữ dữ liệu dưới định dạng JSON/BSON linh hoạt. Thư viện Mongoose đóng vai trò là lớp ánh xạ đối tượng dữ liệu (Object Document Mapper - ODM) giúp định nghĩa Schema chặt chẽ và quản lý các kết nối tối ưu.

\textbf{Giải quyết yêu cầu Chương 2:} 
MongoDB giải quyết trọn vẹn yêu cầu lưu trữ thông tin sản phẩm đa nhà bán hàng với các cấu trúc biến thể phức tạp (màu sắc, kích thước, giá bán riêng biệt). Đặc biệt, MongoDB hỗ trợ các câu lệnh cập nhật nguyên tử ở cấp độ đơn tài liệu (\textit{Single-document Atomicity}) thông qua cú pháp cập nhật có điều kiện như toán tử \texttt{\$inc} kết hợp điều kiện lọc giá trị lớn hơn hoặc bằng lượng mua thực tế ngay dưới tầng lưu trữ:
\begin{verbatim}
db.products.updateOne(
  { _id: productId, "variants.key": variantKey, "variants.stock": { $gte: quantity } },
  { $inc: { "variants.$.stock": -quantity, "totalStock": -quantity } }
)
\end{verbatim}
Câu lệnh này giúp việc kiểm tra tồn kho và trừ kho diễn ra trong một chu kỳ thao tác duy nhất ở tầng dữ liệu, từ đó giảm đáng kể nguy cơ tranh chấp dữ liệu dẫn đến lỗi \textit{Overselling} dưới tải cao.

\textbf{Các phương án công nghệ thay thế và so sánh:}
\begin{itemize}
    \item \textbf{Giải pháp thay thế:} Hệ quản trị cơ sở dữ liệu quan hệ truyền thống (RDBMS) như MySQL hoặc PostgreSQL \cite{rdbms-vs-nosql}.
    \item \textbf{So sánh đối chiếu:}
    \begin{longtable}{|p{3cm}|p{5.5cm}|p{6cm}|}
    \hline
    \textbf{Tiêu chí} & \textbf{RDBMS (MySQL / PostgreSQL)} & \textbf{MongoDB (Đề xuất trong đồ án)} \\ \hline
    \textbf{Mô hình dữ liệu} & Lược đồ tĩnh, yêu cầu liên kết nhiều bảng (\textit{JOIN}) phức tạp để biểu diễn biến thể. & Lược đồ động, lưu trữ biến thể lồng nhau (\textit{Nested Document}) trực quan, không cần JOIN. \\ \hline
    \textbf{Cơ chế khóa tải cao} & Sử dụng cơ chế khóa dòng vật lý (\textit{Pessimistic Locking} qua \texttt{SELECT FOR UPDATE}) gây tắc nghẽn luồng xử lý. & Cơ chế khóa lạc quan không chặn ở cấp độ tài liệu kèm toán tử nguyên tử \texttt{\$inc} cho hiệu năng ghi vượt trội. \\ \hline
    \textbf{Khả năng mở rộng} & Mở rộng theo chiều dọc (\textit{Scaling Up}) tốn kém phần cứng; phân mảnh dữ liệu phức tạp. & Hỗ trợ phân mảnh tự động (\textit{Sharding}) mở rộng theo chiều ngang (\textit{Scaling Out}) dễ dàng. \\ \hline
    \end{longtable}
    \item \textbf{Lý do chọn lựa:} MongoDB không chỉ giúp cấu trúc dữ liệu sản phẩm đa nhà bán hàng trở nên gọn nhẹ, tự nhiên hơn mà còn bảo đảm thông lượng giao dịch ghi ở tải cao lớn hơn nhiều lần so với các hệ cơ sở dữ liệu quan hệ truyền thống nhờ cơ chế cập nhật nguyên tử gọn nhẹ, không gây trễ khóa luồng.
\end{itemize}

\section{Nền tảng lập trình Backend hiệu năng cao và Real-time Communication}
\label{section:3.2}

\subsection{Nền tảng Node.js và Express Framework}
Node.js là một môi trường chạy JavaScript mã nguồn mở hoạt động trên nền bộ máy V8 Engine của Google Chrome. Node.js áp dụng kiến trúc hướng sự kiện (\textit{Event-driven}) và mô hình vào/ra không chặn (\textit{Non-blocking I/O}) hoạt động trên một luồng duy nhất (\textit{Single-thread}) thông qua vòng lặp sự kiện (\textit{Event Loop}).

\textbf{Giải quyết yêu cầu Chương 2:}
Hệ thống sàn đa nhà bán hàng đòi hỏi backend phải xử lý nhiều kết nối đồng thời từ khách hàng và đối tác bán hàng mà vẫn duy trì thời gian phản hồi ổn định. Node.js tận dụng cơ chế bất đồng bộ giúp máy chủ tiếp nhận yêu cầu thanh toán, cập nhật kho, gửi thông báo mà không bị khóa luồng (\textit{Thread blocking}) như một số kiến trúc đa luồng truyền thống. Express framework gọn nhẹ giúp thiết lập hệ thống RESTful API và phân chia định tuyến rõ ràng.

\textbf{Các phương án công nghệ thay thế và so sánh:}
\begin{itemize}
    \item \textbf{Giải pháp thay thế:} Java Spring Boot hoặc Python Django \cite{springboot-vs-nodejs}.
    \item \textbf{So sánh đối chiếu:}
    \begin{itemize}
        \item \textbf{Java Spring Boot:} Là nền tảng mạnh cho hệ thống doanh nghiệp, nhưng thường yêu cầu cấu hình vận hành và bộ nhớ lớn hơn. Mô hình đa luồng truyền thống (\textit{Thread-per-request}) cũng cần được điều chỉnh cẩn thận khi số lượng kết nối đồng thời tăng cao.
        \item \textbf{Python Django:} Hỗ trợ nhiều thư viện mạnh mẽ nhưng hiệu năng xử lý tác vụ I/O đồng thời kém hơn do bị giới hạn bởi khóa GIL (Global Interpreter Lock).
    \end{itemize}
    \item \textbf{Lý do chọn lựa:} Node.js phù hợp với bài toán I/O-intensive của sàn thương mại điện tử, cho phép xây dựng backend thống nhất bằng JavaScript và đáp ứng tốt các yêu cầu kết nối đồng thời, giao tiếp thời gian thực và tích hợp thanh toán.
\end{itemize}

\subsection{Kênh truyền thông thời gian thực Socket.IO}
Socket.IO là một thư viện JavaScript cho phép giao tiếp hai chiều, thời gian thực và dựa trên sự kiện giữa trình duyệt và máy chủ. Nó hoạt động trên nền tảng giao thức WebSocket nhưng bổ sung các cơ chế dự phòng an toàn.

\textbf{Giải quyết yêu cầu Chương 2:}
Socket.IO được sử dụng để hỗ trợ giao tiếp tức thời giữa khách hàng và thương nhân trong phân hệ chat, đồng thời phục vụ các luồng thông báo thời gian thực tại phía client khi backend phát sinh sự kiện mới. Ở tầng kết nối, máy chủ kiểm tra access token được gửi trong phần handshake để gắn socket với người dùng hợp lệ trước khi tham gia phòng chat hoặc nhận sự kiện.

\textbf{Các phương án công nghệ thay thế và so sánh:}
\begin{itemize}
    \item \textbf{Giải pháp thay thế:} HTTP Long Polling hoặc Server-Sent Events (SSE) \cite{websocket-standard}.
    \item \textbf{So sánh đối chiếu:} HTTP Long Polling liên tục gửi yêu cầu kéo (\textit{Pull}) dữ liệu làm nghẽn băng thông mạng và quá tải CPU máy chủ. SSE chỉ hỗ trợ truyền tải dữ liệu một chiều từ Server về Client. Socket.IO thiết lập kết nối TCP toàn song công (\textit{Full-duplex}) duy nhất, giảm thiểu tối đa kích thước header gói tin và hỗ trợ cơ chế tự động kết nối lại (\textit{Auto-reconnection}) thông minh.
    \item \textbf{Lý do chọn lựa:} Socket.IO phù hợp với yêu cầu chat và cập nhật trạng thái theo thời gian thực của hệ thống, đồng thời cho phép tích hợp xác thực access token ở bước handshake và quản lý room thuận tiện cho bài toán nhiều người dùng, nhiều cuộc hội thoại.
\end{itemize}

\section{Lớp đệm bộ nhớ cache và Cơ chế bảo mật đa tầng Redis}
\label{section:3.3}

\subsection{Công nghệ lưu trữ bộ nhớ đệm Redis}
Redis (Remote Dictionary Server) là một hệ thống lưu trữ dữ liệu dạng Key-Value trong bộ nhớ RAM (\textit{In-memory Database}) có mã nguồn mở, hỗ trợ nhiều cấu trúc dữ liệu đa dạng và đạt tốc độ đọc ghi cực cao (dưới $1$ mili-giây).

\textbf{Giải quyết yêu cầu Chương 2:}
Redis được tích hợp tại tầng backend để giải quyết hai nhóm vấn đề cốt lõi:
\begin{itemize}
    \item \textbf{Bảo mật Rate-limiting:} Theo dõi số lần gửi yêu cầu của từng địa chỉ IP hoặc email trong các khoảng thời gian cấu hình sẵn, từ đó giới hạn các API nhạy cảm như đăng nhập, xác minh email, quên mật khẩu, refresh token hoặc truy vấn thống kê vendor.
    \item \textbf{Quản lý phiên xác thực phía server:} Lưu bản ghi phiên theo \texttt{sid} cùng \texttt{currentRefreshJti} để hỗ trợ refresh token rotation, xác minh access token còn gắn với phiên hợp lệ và thu hồi toàn bộ phiên đăng nhập khi người dùng đăng xuất hoặc khi refresh token không còn khớp.
\end{itemize}

\textbf{Các phương án công nghệ thay thế và so sánh:}
\begin{itemize}
    \item \textbf{Giải pháp thay thế:} Lưu trữ bộ nhớ cache trực tiếp trong RAM nội bộ của Node.js (\textit{In-memory Node.js Cache}) hoặc sử dụng Memcached \cite{redis-vs-memcached}.
    \item \textbf{So sánh đối chiếu:} RAM nội bộ của Node.js sẽ bị xóa hoàn toàn khi khởi động lại máy chủ và không thể chia sẻ dữ liệu giữa các tiến trình khi scale-up hệ thống bằng PM2. Memcached tuy nhanh nhưng chỉ hỗ trợ kiểu dữ liệu chuỗi đơn giản, thiếu các cấu trúc dữ liệu nâng cao như Set, Sorted Set của Redis.
    \item \textbf{Lý do chọn lựa:} Redis cung cấp lớp lưu trữ bộ nhớ ngoài có độ trễ thấp, phù hợp cho các thao tác đếm theo cửa sổ thời gian và lưu trạng thái phiên xác thực ngắn hạn. Cách tiếp cận này giúp giảm tải cho MongoDB và tách biệt rõ dữ liệu giao dịch lâu dài với dữ liệu phiên ngắn hạn.
\end{itemize}

\section{Cơ chế khuyến nghị cá nhân hóa lai có suy hao thời gian}
\label{section:3.4}

\subsection{Cơ sở lý thuyết hệ thống gợi ý cá nhân hóa}
Hệ thống khuyến nghị sản phẩm đóng vai trò then chốt trong việc kích cầu tiêu dùng và tăng trưởng doanh số. Hai phương pháp tiếp cận phổ biến nhất là:
\begin{itemize}
    \item \textbf{Lọc dựa trên nội dung (Content-based Filtering):} Gợi ý các sản phẩm có thuộc tính, danh mục tương tự với những sản phẩm người dùng đã tương tác trong quá khứ \cite{recommender-systems}.
    \item \textbf{Lọc cộng tác (Collaborative Filtering):} Tìm kiếm các hành vi tương đồng giữa những người dùng có chung sở thích để đưa ra gợi ý sản phẩm phù hợp.
\end{itemize}

\subsection{Khuyến nghị lai dựa trên tương tác người dùng và suy hao theo thời gian}
Để giảm ảnh hưởng của bài toán khởi đầu lạnh (\textit{Cold Start}) và vẫn tận dụng được lịch sử hành vi người dùng, đồ án sử dụng cơ chế khuyến nghị lai gồm ba bước chính:
\begin{itemize}
    \item \textbf{Bước 1 - Ghi nhận tương tác:} Mỗi lần người dùng xem, nhấp, thêm vào giỏ, thêm yêu thích, đánh giá, mua hàng hoặc dành thời gian trên sản phẩm, hệ thống cập nhật hồ sơ tương tác tương ứng.
    \item \textbf{Bước 2 - Tính điểm tương tác có suy hao thời gian:} Mỗi hồ sơ người dùng - sản phẩm được quy đổi thành một \textit{interactionScore} dựa trên trọng số hành vi và hệ số suy hao theo thời gian.
    \item \textbf{Bước 3 - Sinh danh sách gợi ý lai:} Backend kết hợp hai nguồn gợi ý chính là \textit{item-based collaborative filtering} và \textit{content-based filtering}; nếu số lượng sản phẩm chưa đủ, hệ thống bổ sung thêm các sản phẩm phổ biến đang hoạt động như một lớp fallback.
\end{itemize}

Trong cài đặt hiện tại, điểm tương tác của một cặp người dùng - sản phẩm được tính theo dạng:
\begin{equation}
    Score_{interaction} = Score_{raw} \times e^{- \Delta t / 30}
\end{equation}
Trong đó:
\begin{itemize}
    \item $Score_{raw}$ là tổng điểm thô từ các hành vi như mua hàng, đánh giá, thêm vào giỏ, yêu thích, xem hoặc nhấp chuột.
    \item $\Delta t$ là số ngày tính từ lần tương tác gần nhất.
\end{itemize}

\textbf{Giải quyết yêu cầu Chương 2:}
\begin{itemize}
    \item \textbf{Lọc dựa trên nội dung (Content-based):} Hệ thống lấy các sản phẩm mà người dùng đã tương tác mạnh, sau đó mở rộng sang các sản phẩm cùng danh mục hoặc danh mục con.
    \item \textbf{Lọc cộng tác dựa trên sản phẩm (Item-based Collaborative):} Hệ thống tìm các người dùng có giao cắt hành vi với sản phẩm tương tự, sau đó cộng dồn \textit{interactionScore} để xếp hạng ứng viên.
    \item \textbf{Fallback theo độ phổ biến:} Khi dữ liệu cá nhân hóa chưa đủ, hệ thống bổ sung sản phẩm đang hoạt động theo các tiêu chí như số bán, điểm đánh giá và thời điểm tạo.
\end{itemize}

\textbf{Lý do chọn lựa:} Cách tiếp cận này bám sát dữ liệu hành vi thực tế đang được thu thập trong hệ thống, vừa hỗ trợ cá nhân hóa cho người dùng có lịch sử tương tác, vừa giữ được khả năng fallback ổn định khi dữ liệu còn ít.

\section{Công cụ giao dịch tài chính và Giao diện tương tác nâng cao}
\label{section:3.5}

\subsection{Cổng thanh toán Stripe và VNPay}
Cổng thanh toán điện tử đóng vai trò trung gian bảo mật luồng giao dịch tài chính giữa tài khoản khách hàng, sàn thương mại điện tử và tài khoản ngân hàng của thương nhân.

\textbf{Giải quyết yêu cầu Chương 2:}
Hệ thống tích hợp hai cổng thanh toán lớn: Stripe phục vụ cho các thanh toán thẻ tín dụng quốc tế theo chuẩn PCI-DSS; VNPay cung cấp phương thức thanh toán nội địa thông qua QR và Mobile Banking. Trong cài đặt hiện tại, Stripe xác nhận kết quả bằng webhook ở backend, còn VNPay sử dụng luồng callback hoặc redirect có chữ ký số để backend kiểm tra tính hợp lệ trước khi cập nhật đơn hàng.

\textbf{Phương án thay thế và so sánh:}
\begin{itemize}
    \item \textbf{Giải pháp thay thế:} Chỉ sử dụng hình thức thanh toán thủ công (Chuyển khoản ngân hàng chụp ảnh hóa đơn gửi Admin duyệt).
    \item \textbf{Lý do chọn lựa:} Thanh toán thủ công tốn thời gian, dễ sai lệch dữ liệu và không thuận lợi cho việc giữ kho hoặc hoàn kho tự động theo thời hạn thanh toán. Tích hợp Stripe và VNPay giúp chuẩn hóa luồng thanh toán, đồng bộ kết quả vào backend và hỗ trợ xử lý đơn hàng nhất quán hơn.
\end{itemize}

\subsection{React.js và Thư viện biểu đồ Recharts dành cho Dashboard Merchant}
React.js là một thư viện JavaScript mã nguồn mở dùng để xây dựng giao diện người dùng dựa trên cơ chế thành phần (\textit{Component-based}) và sử dụng cây DOM ảo (\textit{Virtual DOM}) tối ưu hóa tốc độ kết xuất dữ liệu. Recharts là một thư viện vẽ biểu đồ chuyên biệt dành riêng cho hệ sinh thái React, hoạt động dựa trên các thành phần đồ họa vector SVG.

\textbf{Giải quyết yêu cầu Chương 2:}
React.js hỗ trợ xây dựng giao diện Single Page Application với khả năng cập nhật trạng thái linh hoạt và tái sử dụng component cao. Recharts giúp Merchant Console vẽ các biểu đồ doanh thu, thống kê đơn hàng và cảnh báo tồn kho dưới dạng trực quan, đồng thời hỗ trợ hiển thị đáp ứng (\textit{Responsive}) trên nhiều kích thước màn hình.

\textbf{Lý do chọn lựa:} Recharts phù hợp với hệ sinh thái React của đồ án, có cách sử dụng theo component rõ ràng và thuận tiện cho việc xây dựng dashboard bán hàng. Việc dùng SVG cũng giúp biểu đồ hiển thị tốt khi thay đổi kích thước hoặc nhúng vào các layout quản trị khác nhau.

\section*{Kết chương}
Chương 3 đã trình bày cơ sở lý thuyết và các công nghệ chính được áp dụng trong đồ án, bao gồm MongoDB cho lưu trữ dữ liệu tài liệu và cập nhật nguyên tử, Node.js và Express cho backend bất đồng bộ, Redis cho rate-limit và quản lý phiên xác thực, Socket.IO cho giao tiếp thời gian thực, Stripe và VNPay cho thanh toán trực tuyến, cùng cơ chế khuyến nghị sản phẩm dựa trên tương tác người dùng. Từ nền tảng này, chương tiếp theo sẽ chuyển sang kiến trúc phần mềm chi tiết, thiết kế lớp, mô hình dữ liệu và hiện thực hệ thống ở mức cụ thể hơn.

\end{document}
